\chapter{Architecture et implémentation du système}

\section*{Introduction}
\addcontentsline{toc}{section}{Introduction}

Le chapitre précédent a défini le modèle multi-agents, les responsabilités et les limites d'autorité. Le présent chapitre décrit la mise en œuvre concrète des deux systèmes de la Migration Factory : environnement technique, évolution du prototype Java, organisation des dépôts, agents, orchestrateur, backend, persistance, interface, sécurité et difficultés rencontrées.

Les deux déclinaisons ne possèdent pas le même niveau de maturité. Java/Spring Boot constitue la référence la plus avancée, avec un processus complet, une réparation gouvernée, un assistant et un reporting. Angular est en développement actif : les agents principaux sont implémentés et l'intégration du parcours de bout en bout est en cours. Le texte distingue donc les fonctionnalités validées, partielles et planifiées.

% ============================================================
\section{Environnement technique}
% ============================================================

\subsection{Technologies communes}

Avant de distinguer les chaînes propres à Java/Spring Boot et à Angular, il est nécessaire d'identifier le socle technique partagé par les deux déclinaisons. Ce socle regroupe les composants qui assurent le pilotage du workflow, la persistance de l'état, l'exposition des services et l'interaction avec le développeur. Le tableau suivant présente ces technologies communes et précise leur rôle dans l'architecture globale.

\begin{table}[H]
    \centering
    \small
    \arrayrulecolor{cgigray}
    \rowcolors{2}{cgilight!55}{white}
    \begin{tabularx}{\textwidth}{|>{\bfseries\RaggedRight\arraybackslash}p{3.7cm}|>{\RaggedRight\arraybackslash}p{4.3cm}|>{\RaggedRight\arraybackslash}X|}
        \hline
        \rowcolor{cgilight!95}
        \textbf{Couche} & \textbf{Technologies} & \textbf{Rôle} \\
        \hline
        Interface & React et Next.js & Configuration, cockpit, décisions, assistant et rapports. \\
        \hline
        Backend & Python et FastAPI & API HTTP, schémas, services applicatifs et projections. \\
        \hline
        Orchestration & LangGraph & Graphe d'exécution, interruptions, reprises et routage. \\
        \hline
        Persistance & SQLite et artefacts sur disque & États, événements, décisions, commandes, journaux et rapports. \\
        \hline
        Intelligence & Modèles et fournisseurs configurés selon le dépôt et le rôle & Analyse, planification, révision, réparation et assistance. \\
        \hline
        Collaboration & GitHub, Jira, Teams et Confluence & Versionnement, backlog, échanges et documentation. \\
        \hline
    \end{tabularx}
    \caption{Technologies communes aux deux systèmes}
    \label{tab:technologies-communes-ch4}
    \rowcolors{2}{white}{white}
    \arrayrulecolor{black}
\end{table}

Le rapport ne présente pas Azure OpenAI comme fournisseur commun aux deux dépôts. L'audit Java a confirmé son utilisation dans le runtime observé ; la configuration Angular reste décrite de manière générique tant que son fournisseur courant n'est pas établi par une preuve de code récente.

FastAPI fournit les routes, la validation et l'injection de dépendances \cite{fastapi-dependencies}. LangGraph coordonne les étapes avec état et checkpoints \cite{langgraph-persistence}. SQLite reste adapté au périmètre local ; le mode WAL réduit certains blocages entre lecteurs et écrivain \cite{sqlite-wal}.

\subsection{Environnement Java/Spring Boot}

La chaîne Java mobilise JDK 11, 17 et 21, Maven, des profils Spring Boot et OpenRewrite. Les recettes OpenRewrite sont composables et doivent s'abstenir lorsqu'elles ne peuvent pas établir la sûreté d'une transformation \cite{openrewrite-recipes,openrewrite-best-practices}. Plusieurs scripts PowerShell facilitent le lancement et le diagnostic sous Windows.

La source Java est fournie par un chemin local. Le système vérifie le projet, les JDK, Maven, le profil de migration et le répertoire de sortie avant la création de l'exécution.

\subsection{Environnement Angular}

La chaîne Angular mobilise Node.js, npm, Angular CLI, TypeScript et RxJS. La matrice officielle Angular relie chaque version à des plages compatibles de Node.js, TypeScript et RxJS \cite{angular-versions}. L'installation, la construction et les tests utilisent les cibles du projet \cite{angular-cli-build,angular-cli-test}.

L'entrée prévue est un paquet ZIP contenant l'application source. Le backend crée un snapshot en lecture seule, puis prépare un espace de stage en excluant \texttt{node\_modules}, \texttt{dist}, les caches, les couvertures et les sorties régénérables.

% ============================================================
\section{Évolution technique de la solution Java/Spring Boot}
% ============================================================

L'architecture actuelle résulte de plusieurs incréments successifs. Cette évolution permet de relier les choix techniques aux limites constatées lors des démonstrations.

\subsection{Première version : preuve de concept en ligne de commande}

La première version opérationnelle était principalement utilisée en ligne de commande. Elle associait des scripts de lancement, des agents d'analyse et de planification, des transformations, l'exécution de Maven, l'analyse des tests et la génération d'artefacts. Cette version a validé le principe d'une migration progressive exécutée dans un espace isolé.

La démonstration du 25 mai 2026 a confirmé la faisabilité technique du workflow. Ses limites étaient surtout liées au pilotage : l'utilisateur devait suivre plusieurs commandes, parcourir les artefacts et reconstruire mentalement l'état global de la migration.

\subsection{Consolidation du prototype CLI}

Entre le 25 mai et le 3 juin, le workflow a été corrigé et stabilisé. Les sorties ont été enrichies, le scénario de démonstration a été mieux maîtrisé et plusieurs contrôles ont été renforcés. La deuxième démonstration, réalisée le 3 juin devant une audience CGI élargie incluant davantage de responsables et directeurs, a présenté cette version consolidée.

Cette étape n'a pas encore introduit la Control Tower complète. Elle a surtout validé la robustesse croissante du prototype et confirmé l'intérêt de construire une interface de pilotage.

\subsection{Passage à la Control Tower full stack}

La troisième étape a transformé la preuve de concept en plateforme. Un backend FastAPI a centralisé les API, les transitions, la persistance et les projections. SQLite a conservé les migrations, stages, commandes, événements et décisions. Un frontend React/Next.js a fourni un cockpit, tandis qu'un chatbot contextualisé a permis d'expliquer l'état courant.

La démonstration du 19 juin 2026 a présenté cette Control Tower full stack : backend, frontend, persistance, pipeline, points de décision, journaux et assistant. Cette version a déplacé le projet d'une logique d'outil technique vers une logique de plateforme de migration.

\subsection{Consolidation de la V1 gouvernée}

Les travaux engagés à partir du 20 juin préparent une V1 prévue pour septembre 2026. Ils portent sur :

\begin{itemize}
    \item un nouveau frontend améliorant la visualisation du workflow et des preuves ;
    \item une boucle de réparation séparant le proposeur et l'agent de révision ;
    \item une approbation humaine liée à l'empreinte exacte du correctif ;
    \item une application déterministe dans un espace isolé ;
    \item une nouvelle compilation et l'exécution des tests après chaque correctif ;
    \item une amélioration des prompts, du contexte, des sorties structurées et de l'exploitation des LLM ;
    \item un reporting enrichi et une meilleure observation de la consommation des modèles.
\end{itemize}

Cette quatrième étape est encore en cours à la date de référence du rapport ; elle ne doit donc pas être présentée comme une démonstration déjà réalisée.

\begin{figure}[H]
    \centering
    \resizebox{0.98\textwidth}{!}{%
    \begin{tikzpicture}[
        box/.style={draw=cgiblue,rounded corners=3pt,fill=cgilight!55,text width=3.2cm,minimum height=1.25cm,align=center,font=\scriptsize\bfseries},
        future/.style={draw=cgired,rounded corners=3pt,fill=white,text width=3.2cm,minimum height=1.25cm,align=center,font=\scriptsize\bfseries},
        arr/.style={-{Latex[length=2mm]},thick,draw=cgigray}
    ]
        \node[box] (cli1) {Prototype CLI\\analyse, plan, Maven};
        \node[box,right=0.35cm of cli1] (cli2) {CLI consolidée\\contrôles et sorties};
        \node[box,right=0.35cm of cli2] (api) {Backend FastAPI\\SQLite et événements};
        \node[box,right=0.35cm of api] (ui) {Control Tower\\frontend et chatbot};
        \node[future,right=0.35cm of ui] (v1) {V1 gouvernée\\nouveau frontend et réparation};
        \draw[arr] (cli1)--(cli2);
        \draw[arr] (cli2)--(api);
        \draw[arr] (api)--(ui);
        \draw[arr] (ui)--(v1);
    \end{tikzpicture}%
    }
    \caption{Évolution technique de la solution Java/Spring Boot}
    \label{fig:evolution-technique-java-ch4}
\end{figure}

\scriptsize
\begin{table}[H]
    \centering
    \arrayrulecolor{cgigray}
    \begin{tabularx}{\textwidth}{|>{\bfseries\RaggedRight\arraybackslash}p{2.4cm}|>{\RaggedRight\arraybackslash}p{3cm}|>{\RaggedRight\arraybackslash}p{3cm}|>{\RaggedRight\arraybackslash}p{3cm}|>{\RaggedRight\arraybackslash}X|}
        \hline
        \rowcolor{cgilight!95}
        \textbf{Dimension} & \textbf{J1 -- 25 mai} & \textbf{J2 -- 3 juin} & \textbf{J3 -- 19 juin} & \textbf{J4 -- septembre} \\
        \hline
        Interface & Ligne de commande & CLI améliorée & Frontend web & Nouveau frontend \\
        \hline
        Backend & Workflow local & Workflow stabilisé & FastAPI et SQLite & Services consolidés \\
        \hline
        Assistant & Non intégré & Assistance limitée & Chatbot contextualisé & Assistant enrichi \\
        \hline
        Réparation & Correction manuelle & Diagnostic amélioré & Premiers composants & Boucle gouvernée complète \\
        \hline
        LLM & Premiers agents & Agents consolidés & Agents intégrés à la plateforme & Contexte et sorties renforcés \\
        \hline
        Validation & Maven et tests & Scénario consolidé & Pilotage full stack & Revalidation après correctif \\
        \hline
        Statut & Réalisé & Réalisé & Réalisé & Planifié \\
        \hline
    \end{tabularx}
    \caption{Comparaison des principaux incréments Java/Spring Boot}
    \label{tab:comparaison-demos-java-ch4}
    \arrayrulecolor{black}
\end{table}
\normalsize

% ============================================================
\section{Structure simplifiée des dépôts}
% ============================================================

Les arbres suivants représentent une vue simplifiée, et non la liste exhaustive de tous les dossiers.

\subsection{Dépôt Java/Spring Boot}

\begin{verbatim}
migration-factory-java/
|-- migration_factory/
|   |-- control_tower/
|   |   |-- adapters/fastapi/
|   |   |-- application/
|   |   |-- domain/
|   |   `-- infrastructure/sqlite/
|   |-- orchestrator/
|   `-- agents/
|-- web/control-tower/
|-- tests/control_tower/
|-- scripts/
`-- docs/
\end{verbatim}

Le domaine et les services applicatifs sont séparés des adaptateurs HTTP et SQLite. Les agents sont isolés des composants d'exécution et ne pilotent pas directement le frontend.

\subsection{Dépôt Angular}

\begin{verbatim}
angular-migration/
|-- backend/
|   |-- app/
|   |   |-- domain/
|   |   |-- services/
|   |   |-- api/
|   |   |-- infrastructure/
|   |   `-- orchestration/
|   `-- tests/
|-- frontend/
|-- goals/
|-- artifacts/
|-- scripts/
`-- tests/
\end{verbatim}

Les deux dépôts partagent des principes et une terminologie, mais pas leurs services d'exécution. Cette séparation réduit le couplage entre Maven/JDK et npm/Node.js.

\subsection{Déploiement local}

Les deux sous-sections précédentes ont présenté l'organisation interne des dépôts Java/Spring Boot et Angular. La présente vue complète cette description en montrant comment leurs composants sont réellement instanciés et reliés pendant l'exécution locale : le navigateur accède au frontend, le frontend dialogue avec le backend, et le backend coordonne persistance, espaces isolés, outils d'écosystème et fournisseur LLM. Elle fait ainsi le lien entre la structure du code et l'architecture d'exécution.

\begin{figure}[H]
    \centering
    \resizebox{0.94\textwidth}{!}{%
    \begin{tikzpicture}[
        node/.style={draw=cgiblue,rounded corners=3pt,fill=cgilight!55,minimum width=4cm,minimum height=1.15cm,align=center,font=\small\bfseries},
        store/.style={draw=cgigray,rounded corners=3pt,fill=white,minimum width=4cm,minimum height=1.15cm,align=center,font=\small\bfseries},
        arr/.style={-{Latex[length=2mm]},thick,draw=cgigray}
    ]
        \node[node] (browser) {Navigateur\\cockpit local};
        \node[node,right=1cm of browser] (frontend) {Serveur frontend\\Next.js};
        \node[node,right=1cm of frontend] (backend) {Backend local\\FastAPI};
        \node[store,below=1.1cm of backend] (db) {SQLite et artefacts};
        \node[store,left=1cm of db] (workspace) {Espaces isolés\\et outils};
        \node[store,right=1cm of db] (llm) {Fournisseur LLM\\externe};
        \draw[arr] (browser)--(frontend);
        \draw[arr] (frontend)--(backend);
        \draw[arr] (backend)--(db);
        \draw[arr] (backend)--(workspace);
        \draw[arr] (backend)--(llm);
    \end{tikzpicture}%
    }
    \caption{Déploiement local des prototypes}
    \label{fig:deploiement-local-ch4}
\end{figure}

% ============================================================
\section{Implémentation des agents}
% ============================================================

Chaque véritable agent est présenté selon le même format : objectif, entrées, traitement, sorties, scénario, autorité et limitations.

\subsection{Agent d'analyse}

\begin{table}[H]
    \centering
    \small
    \begin{tabularx}{\textwidth}{|>{\bfseries\RaggedRight\arraybackslash}p{3.1cm}|>{\RaggedRight\arraybackslash}X|}
        \hline
        Objectif & Construire une représentation exploitable de l'application sans modifier la source. \\
        \hline
        Entrées & Snapshot, versions, manifestes, structure, configurations et profil. \\
        \hline
        Traitement & Inventaire déterministe, sélection d'extraits, raisonnement contextualisé et validation de schéma. \\
        \hline
        Sorties & Rapport d'analyse, risques, références et preuves. \\
        \hline
        Exemple & Détecter une version Spring Boot ancienne ou une incompatibilité Angular--Node.js. \\
        \hline
        Autorité & Lecture uniquement ; aucune commande ni écriture sur la source. \\
        \hline
        Limitations & Dépendances privées et configurations dynamiques peuvent nécessiter une information humaine. \\
        \hline
    \end{tabularx}
    \caption{Fiche de l'agent d'analyse}
    \label{tab:agent-analyse-ch4}
\end{table}

\subsection{Agent de révision de l'analyse}

\begin{table}[H]
    \centering
    \small
    \begin{tabularx}{\textwidth}{|>{\bfseries\RaggedRight\arraybackslash}p{3.1cm}|>{\RaggedRight\arraybackslash}X|}
        \hline
        Objectif & Vérifier couverture, cohérence et exploitabilité du rapport. \\
        \hline
        Entrées & Analyse, preuves, profil et règles d'acceptation. \\
        \hline
        Traitement & Comparaison des affirmations avec les preuves et détection des omissions. \\
        \hline
        Sorties & Acceptation, demande de révision ou rejet motivé. \\
        \hline
        Exemple & Refuser un rapport Angular qui ne relie pas Angular à Node.js et TypeScript. \\
        \hline
        Autorité & Ne lance aucune commande et ne modifie pas l'application. \\
        \hline
        Limitations & Sa qualité dépend des preuves disponibles et du schéma d'analyse. \\
        \hline
    \end{tabularx}
    \caption{Fiche de l'agent de révision de l'analyse}
    \label{tab:agent-review-analyse-ch4}
\end{table}

\subsection{Agent de planification}

\begin{table}[H]
    \centering
    \small
    \begin{tabularx}{\textwidth}{|>{\bfseries\RaggedRight\arraybackslash}p{3.1cm}|>{\RaggedRight\arraybackslash}X|}
        \hline
        Objectif & Décomposer la cible en stages, unités de travail et validations. \\
        \hline
        Entrées & Analyse validée, cible, catalogue de compatibilité et politiques. \\
        \hline
        Traitement & Sélection d'une trajectoire, ordonnancement et définition des critères de sortie. \\
        \hline
        Sorties & Plan versionné, dépendances, risques et points de validation. \\
        \hline
        Exemple & Construire Spring Boot 2.1 $\rightarrow$ 2.7 $\rightarrow$ 3.5 $\rightarrow$ 4.0. \\
        \hline
        Autorité & Ne transforme pas le projet ; le plan doit être révisé et approuvé. \\
        \hline
        Limitations & Dépend de la qualité de l'analyse et des profils disponibles. \\
        \hline
    \end{tabularx}
    \caption{Fiche de l'agent de planification}
    \label{tab:agent-planification-ch4}
\end{table}

\subsection{Agent de révision du plan}

\begin{table}[H]
    \centering
    \small
    \begin{tabularx}{\textwidth}{|>{\bfseries\RaggedRight\arraybackslash}p{3.1cm}|>{\RaggedRight\arraybackslash}X|}
        \hline
        Objectif & Contrôler l'ordre, les dépendances et les validations. \\
        \hline
        Entrées & Plan, analyse, trajectoire et contraintes. \\
        \hline
        Traitement & Vérification des prérequis et de la couverture des risques. \\
        \hline
        Sorties & Acceptation ou demande de nouvelle version. \\
        \hline
        Exemple & Refuser un saut Angular majeur non autorisé par la trajectoire. \\
        \hline
        Autorité & Aucune écriture sur le projet et aucune approbation humaine implicite. \\
        \hline
        Limitations & Ne peut pas prouver la réussite future d'une compilation. \\
        \hline
    \end{tabularx}
    \caption{Fiche de l'agent de révision du plan}
    \label{tab:agent-review-plan-ch4}
\end{table}

\subsection{Agent de transformation}

\begin{table}[H]
    \centering
    \small
    \begin{tabularx}{\textwidth}{|>{\bfseries\RaggedRight\arraybackslash}p{3.1cm}|>{\RaggedRight\arraybackslash}X|}
        \hline
        Objectif & Produire les changements candidats décrits par un plan approuvé. \\
        \hline
        Entrées & Plan, espace isolé, profil et règles de transformation. \\
        \hline
        Traitement & Transformations déterministes et génération de changements contextualisés. \\
        \hline
        Sorties & Fichiers modifiés, diff, métadonnées et preuves. \\
        \hline
        Exemple & Adapter dépendances, configuration Spring ou code TypeScript. \\
        \hline
        Autorité & Écriture limitée à l'espace isolé fourni par le backend. \\
        \hline
        Limitations & Les modifications restent candidates tant que compilation et tests ne réussissent pas. \\
        \hline
    \end{tabularx}
    \caption{Fiche de l'agent de transformation}
    \label{tab:agent-transformation-ch4}
\end{table}

\subsection{Proposeur de réparation}

\begin{table}[H]
    \centering
    \small
    \begin{tabularx}{\textwidth}{|>{\bfseries\RaggedRight\arraybackslash}p{3.1cm}|>{\RaggedRight\arraybackslash}X|}
        \hline
        Objectif & Générer un correctif candidat à partir d'un échec vérifié. \\
        \hline
        Entrées & Erreurs, commande, extraits, profil et état de l'espace isolé. \\
        \hline
        Traitement & Diagnostic, sélection des fichiers et production d'un diff structuré. \\
        \hline
        Sorties & Diff candidat, justification et fichiers concernés. \\
        \hline
        Exemple & Corriger une incompatibilité de compilation Java. \\
        \hline
        Autorité & Ne peut ni approuver ni appliquer. \\
        \hline
        Limitations & Le contexte borné peut omettre une dépendance indirecte. \\
        \hline
    \end{tabularx}
    \caption{Fiche du proposeur de réparation}
    \label{tab:agent-proposeur-reparation-ch4}
\end{table}

\subsection{Agent de révision de la réparation}

\begin{table}[H]
    \centering
    \small
    \begin{tabularx}{\textwidth}{|>{\bfseries\RaggedRight\arraybackslash}p{3.1cm}|>{\RaggedRight\arraybackslash}X|}
        \hline
        Objectif & Évaluer la proposition et produire, si nécessaire, un diff consolidé. \\
        \hline
        Entrées & Proposition, erreur, politiques et preuves. \\
        \hline
        Traitement & Contrôle de cohérence, risque, portée et applicabilité. \\
        \hline
        Sorties & Acceptation, rejet ou diff consolidé. \\
        \hline
        Exemple & Réduire un correctif trop large à la cause de l'échec. \\
        \hline
        Autorité & Ne peut pas appliquer ; l'humain décide sur la version exacte. \\
        \hline
        Limitations & Une sortie correcte en apparence peut encore échouer à la compilation. \\
        \hline
    \end{tabularx}
    \caption{Fiche de l'agent de révision de la réparation}
    \label{tab:agent-review-reparation-ch4}
\end{table}

\subsection{Assistant de migration}

\begin{table}[H]
    \centering
    \small
    \begin{tabularx}{\textwidth}{|>{\bfseries\RaggedRight\arraybackslash}p{3.1cm}|>{\RaggedRight\arraybackslash}X|}
        \hline
        Objectif & Expliquer l'état, les blocages et les actions possibles. \\
        \hline
        Entrées & Projection, événements, preuves autorisées et historique borné. \\
        \hline
        Traitement & Synthèse et réponse contextualisée en lecture seule. \\
        \hline
        Sorties & Explication, résumé ou recommandation d'action. \\
        \hline
        Exemple & Expliquer pourquoi une migration attend une approbation. \\
        \hline
        Autorité & Ne peut ni exécuter, ni approuver, ni modifier des fichiers. \\
        \hline
        Limitations & Une information absente de l'état ne doit pas être inventée. \\
        \hline
    \end{tabularx}
    \caption{Fiche de l'assistant de migration}
    \label{tab:agent-assistant-ch4}
\end{table}

% ============================================================
\section{Services déterministes et orchestrateur}
% ============================================================

Le résolveur de profils, le service de transitions, l'exécuteur, les services de compilation et de tests, le stockage et le reporting ne sont pas des agents. Ils implémentent des règles vérifiables.

Le frontend transmet une intention à l'API. Le backend valide le schéma, charge l'état et vérifie l'action. L'orchestrateur choisit ensuite le nœud autorisé. Les résultats sont persistés avant toute projection vers le cockpit.

\begin{figure}[H]
    \centering
    \resizebox{\textwidth}{!}{%
    \begin{tikzpicture}[
        step/.style={draw=cgiblue,rounded corners=3pt,fill=cgilight!55,minimum width=2.75cm,minimum height=1cm,align=center,font=\small\bfseries},
        gate/.style={draw=cgired,rounded corners=3pt,fill=white,minimum width=2.75cm,minimum height=1cm,align=center,font=\small\bfseries},
        arr/.style={-{Latex[length=2mm]},thick,draw=cgigray}
    ]
        \node[step] (qualify) {Qualification};
        \node[step,right=0.4cm of qualify] (analyse) {Analyse};
        \node[step,right=0.4cm of analyse] (review) {Révision};
        \node[step,right=0.4cm of review] (plan) {Planification};
        \node[gate,right=0.4cm of plan] (approval) {Décision};
        \node[step,right=0.4cm of approval] (transform) {Transformation};
        \node[step,right=0.4cm of transform] (validate) {Validation};
        \draw[arr] (qualify)--(analyse);
        \draw[arr] (analyse)--(review);
        \draw[arr] (review)--(plan);
        \draw[arr] (plan)--(approval);
        \draw[arr] (approval)--(transform);
        \draw[arr] (transform)--(validate);
    \end{tikzpicture}%
    }
    \caption{Séquencement principal implémenté par l'orchestrateur}
    \label{fig:sequencement-orchestrateur-ch4}
\end{figure}

\subsection{Contrats techniques simplifiés}

La représentation suivante illustre le contrat logique d'une décision. Les noms sont simplifiés pour le rapport ; ils ne constituent pas une copie exhaustive d'une route.

\begin{verbatim}
{
  "run_id": "...",
  "gate_id": "...",
  "decision": "approve | reject | revise",
  "expected_checksum": "sha256:...",
  "idempotency_key": "..."
}
\end{verbatim}

Un manifeste de commande contient de manière analogue l'exécutable, les arguments validés, le répertoire de travail, l'environnement matérialisé, le délai maximal et les identifiants de corrélation. L'interface ou un LLM ne transmet jamais une commande shell libre.

% ============================================================
\section{Backend, persistance et preuves}
% ============================================================

Les API couvrent configurations, prérequis, migrations, stages, décisions, réparations, événements, artefacts, assistant et rapports. Les schémas typés détectent les sorties incomplètes.

Les entités principales sont Migration, Stage, Commande, Événement, Artefact, Point de validation, Décision, Proposition, Révision et Rapport. Les événements sont séquencés ; les artefacts conservent type, producteur, chemin et empreinte ; les commandes conservent état, durée, code de sortie et journaux.

\begin{verbatim}
run_<id>/
|-- source_snapshot/
|-- stages/
|   |-- stage_01/workspace/
|   |-- stage_01/analysis/
|   |-- stage_01/planning/
|   `-- stage_01/validation/
|-- repairs/
|-- logs/
|-- evidence/
`-- reports/
\end{verbatim}

La source externe n'est pas utilisée comme répertoire d'écriture. Le candidat validé d'un stage devient la base du suivant. Les décisions sensibles sont liées par SHA-256 à l'artefact examiné.

% ============================================================
\section{Interface, cockpit et captures}
% ============================================================

Le cockpit présente le résumé, le pipeline, le stage sélectionné, l'analyse, le plan, les points de validation, les commandes, les journaux, les événements, la réparation, l'assistant et les rapports. Le frontend projette l'état fourni par le backend ; il ne maintient pas une vérité parallèle.

\begin{figure}[H]
    \centering
    \IfFileExists{figures/chapitre4/java-demo1-cli.png}{%
        \includegraphics[width=0.96\textwidth]{figures/chapitre4/java-demo1-cli.png}%
    }{%
        \fbox{\parbox[c][4.6cm][c]{0.9\textwidth}{\centering\bfseries Capture historique à intégrer : Java J1\\\normalfont Première démonstration de la solution CLI du 25 mai 2026.}}%
    }
    \caption{Première version Java/Spring Boot utilisée en ligne de commande}
    \label{fig:capture-java-demo1-ch4}
\end{figure}

\begin{figure}[H]
    \centering
    \IfFileExists{figures/chapitre4/java-demo3-control-tower.png}{%
        \includegraphics[width=0.96\textwidth]{figures/chapitre4/java-demo3-control-tower.png}%
    }{%
        \fbox{\parbox[c][5cm][c]{0.9\textwidth}{\centering\bfseries Capture réelle à intégrer : Java J3\\\normalfont Control Tower full stack présentée le 19 juin 2026.}}%
    }
    \caption{Control Tower Java/Spring Boot}
    \label{fig:capture-control-tower-java-ch4}
\end{figure}

\begin{figure}[H]
    \centering
    \IfFileExists{figures/chapitre4/java-demo3-chatbot.png}{%
        \includegraphics[width=0.96\textwidth]{figures/chapitre4/java-demo3-chatbot.png}%
    }{%
        \fbox{\parbox[c][5cm][c]{0.9\textwidth}{\centering\bfseries Capture réelle à intégrer : chatbot Java\\\normalfont Assistant contextualisé présenté dans la plateforme full stack.}}%
    }
    \caption{Assistant conversationnel de la Control Tower}
    \label{fig:capture-chatbot-java-ch4}
\end{figure}

\begin{figure}[H]
    \centering
    \IfFileExists{figures/chapitre4/java-v1-reparation.png}{%
        \includegraphics[width=0.96\textwidth]{figures/chapitre4/java-v1-reparation.png}%
    }{%
        \fbox{\parbox[c][5cm][c]{0.9\textwidth}{\centering\bfseries Capture future à intégrer : réparation gouvernée\\\normalfont Proposition, révision, décision et résultat de revalidation de la V1.}}%
    }
    \caption{Écran de réparation gouvernée de la V1 Java}
    \label{fig:capture-reparation-java-ch4}
\end{figure}

\begin{figure}[H]
    \centering
    \IfFileExists{figures/chapitre4/angular-demo-a1.png}{%
        \includegraphics[width=0.96\textwidth]{figures/chapitre4/angular-demo-a1.png}%
    }{%
        \fbox{\parbox[c][5cm][c]{0.9\textwidth}{\centering\bfseries Capture planifiée : démonstration Angular A1\\\normalfont Première version Angular 18--21 prévue le 11 août 2026.}}%
    }
    \caption{Première démonstration d'avancement Angular}
    \label{fig:capture-angular-a1-ch4}
\end{figure}

Ces emplacements rendent explicites les preuves visuelles encore à fournir. Aucune capture fictive n'est présentée comme résultat.

% ============================================================
\section{Authentification, sécurité et limites d'exploitation}
% ============================================================

Les prototypes ne comportent pas d'authentification, de comptes ou de contrôle d'accès par rôle. Le système Java limite l'API au poste local, contrôle l'origine de requêtes mutatives et utilise un identifiant client statique ; ces mécanismes ne constituent pas une identité fiable.

Une industrialisation devra ajouter authentification forte, rôles, autorisation par action, sessions, révocation et journal d'audit attribuable. Ces évolutions sont cohérentes avec les principes de gestion de la sécurité d'ISO/IEC 27001 \cite{iso27001}.

Les contextes LLM sont bornés et nettoyés. La rédaction demeure fondée sur des règles et ne garantit pas la détection de toute valeur sensible intégrée arbitrairement dans le code.

% ============================================================
\section{Difficultés techniques et corrections}
% ============================================================

\begin{table}[H]
    \centering
    \small
    \arrayrulecolor{cgigray}
    \rowcolors{2}{cgilight!55}{white}
    \begin{tabularx}{\textwidth}{|>{\bfseries\RaggedRight\arraybackslash}p{3.6cm}|>{\RaggedRight\arraybackslash}p{4.5cm}|>{\RaggedRight\arraybackslash}X|}
        \hline
        \rowcolor{cgilight!95}
        \textbf{Difficulté} & \textbf{Cause} & \textbf{Correction} \\
        \hline
        Correctif non applicable & Fins de ligne, hunk ou état obsolète. & Diff structuré, normalisation, empreinte et contrôle d'applicabilité. \\
        \hline
        Autorité de commande ambiguë & Plusieurs composants pouvaient influencer l'exécution. & Manifeste autoritaire construit par le backend. \\
        \hline
        Profil Angular rejeté & Catalogue initial incompatible avec Node 22.23.1/npm 10.9.8. & Catalogue versionné et paire explicitement validée. \\
        \hline
        Base SQLite incohérente & Plusieurs têtes de migration ou mauvaise base active. & Identification de la base autoritaire et réconciliation Alembic. \\
        \hline
        Intégration frontend & Fusion de nouvelles vues et fonctions existantes. & Intégration progressive et comparaison au baseline. \\
        \hline
    \end{tabularx}
    \caption{Principales difficultés et corrections}
    \label{tab:difficultes-corrections-ch4}
    \rowcolors{2}{white}{white}
    \arrayrulecolor{black}
\end{table}

% ============================================================
\section{Preuves techniques disponibles}
% ============================================================

Une campagne d'intégration Java liée au nouveau frontend et à l'assistant a produit 607 tests backend réussis et 4 ignorés. La vérification statique des types, l'analyse de conformité et la construction du frontend ont réussi. Les tests frontend comportaient 12 échecs au baseline et 11 après intégration ; cette comparaison a été utilisée pour distinguer les défauts préexistants d'une régression introduite.

Ces valeurs ne représentent pas un benchmark global de la Migration Factory. Elles constituent une preuve ponctuelle d'intégration. Les scénarios Maven réels restent principalement exécutés manuellement ou en recette utilisateur.

Pour Angular, une migration manuelle d'Angular 18 vers 20 a été réalisée afin de comprendre les transformations. Les agents d'analyse, de révision, de planification et de transformation sont développés. Le processus complet et ses métriques quantitatives restent en cours de validation.

Aucun des deux prototypes ne dispose actuellement d'une chaîne CI commune dans le dépôt du produit. Il n'existe pas non plus de politique globale unique de lint couvrant toutes les couches. Ces limites sont reprises au Chapitre V.

\section*{Conclusion}
\addcontentsline{toc}{section}{Conclusion}

Ce chapitre a décrit l'implémentation en distinguant les agents intelligents des services déterministes. Il a également montré que l'architecture actuelle résulte d'une évolution progressive : prototype CLI, consolidation, passage à la Control Tower full stack et préparation d'une V1 gouvernée.

Les fiches uniformes précisent les entrées, traitements, sorties, exemples, autorités et limites. Les contrats simplifiés, les preuves d'intégration et les difficultés techniques renforcent le caractère concret de la réalisation. Le chapitre suivant formalise les démonstrations réalisées ou planifiées et évalue les résultats disponibles, la qualité du système, les limites de mesure et les perspectives.
